\documentclass[11pt]{article}
\usepackage[margin=1in]{geometry}
\usepackage{amsmath,amssymb,amsthm}
\usepackage[T1]{fontenc}
\usepackage{lmodern}
\usepackage{hyperref}
\title{Conjecture 00000000147 Does Not Specify a Probability Model}
\author{quantumtool}
\date{October 19, 2026}
\begin{document}
\maketitle
\section*{The statement is undefined as written}
The conjecture writes $\mathbb P(g_n>t\log p_n)$ but does not specify a probability space, random variable, or sampling rule. Prime gaps $g_n$ and primes $p_n$ are deterministic. Thus the probability in the displayed expression is not defined by the data of the conjecture.

Even if one assigns the deterministic event its natural point-mass probability for each $n$, its value is either 0 or 1. For every fixed $t>0$, however, $e^{-t}$ is strictly between 0 and 1. A sequence of values confined to $\{0,1\}$ cannot converge to this interior value: each term stays at least $\min(e^{-t},1-e^{-t})$ away from it. A meaningful tail conjecture would need to specify a random model or an empirical distribution over a growing range of indices; these choices define different statements.

\section*{Formalization scope}
The Lean file checks the arithmetic core of the point-mass interpretation: a deterministic event probability, encoded as 0 or 1, cannot lie in the central interval corresponding to a fixed positive tolerance around $1/2$. The absence of a probability space and the interpretation choice are discussed above, not represented as a Lean measure-theory theorem.
\section*{Lean formalization}
Lean formalizes the separation of Boolean probabilities (0 or 1) from the open interval around 1/2. The problem statement’s missing sample space and the conditional conclusion about its intended interpretation are explained in the report.
\end{document}
